\begin{textsample}{POS Dim 1 – ms – Score 47.00 – ms/ms\_em\_52.txt}  \label{ex:f1_pos_186}
1.6.
Objetivos Aplicar os conhecimentos das Ciências da Natureza para reconhecer e analisar de forma \textbf{crítica} os processos que abordam : - as principais técnicas de manipulação e \textbf{edição} genética da atualidade e suas aplicações ; - as implicações éticas das principais técnicas de \textbf{edição} genética e o impacto de sua utilização para o desenvolvimento das futuras sociedades ; - a manipulação genética no processo de evolução, tendo como base a \textbf{teoria} \textbf{moderna} da evolução ; - os benefícios e contribuições da manipulação genética para a sociedade, em áreas como a engenharia genética de alimentos, biomedicina, nanotecnologia, dentre outras ; - a utilização da genética e da seleção artificial para justificar os processos de eugenia, discriminação, segregação e superioridade racial ; - a participação em \textbf{debates} e decisões coletivas sobre as vantagens e os limites da \textbf{edição} genética em aplicações de intervenção sociocultural. 1.7.
Relação com outras unidades Ver lista de Unidades Curriculares do Catálogo que se correlacionam a esta Unidade. 1.8.
Perfil docente - Possuir licenciatura em Biologia ; - Ter desenvolvido e/ou desenvolver atividades de estudo e/ou pesquisa no campo de conhecimento correspondente à unidade curricular ( genética humana, biomedicina, bioética, biotecnologia, engenharia genética, ecologia, dentre outros ) ; - Experiências e/ou interesse do professor no campo de pesquisa ; - Conhecimento e/ou disposição para o uso de metodologias ativas e Tecnologias \textbf{Digitais} de Informação e Comunicação ( TDIC ). 1.9.
Recursos - Dispositivos com acesso à internet e \textbf{edição} de texto, lousa \textbf{digital} ou datashow. - Acervo \textbf{impresso} e \textbf{digital} de material de pesquisa. - Parcerias associadas com a sugestão didática : - Participação das atividades de grupo de pesquisa da USP - Bioética e Biotecnologias : uma \textbf{abordagem} \textbf{multidisciplinar} ( www.direitorp.usp.br/wp-content/uploads/2019/08/Grupo-de-Pesquisa-Bioética.pdf ) ; - Laboratórios de Genética e Evolução - INBIO/UFMS.
II. \textbf{Organizador} curricular 2.3.
Sugestões didáticas Sugestões didáticas associadas aos eixos estruturantes Estudo de caso, \textbf{supondo} as seguintes etapas : - \textbf{Análise} do filme Gattaca e de outros correlacionados ( Equilibrium/X-Men - Confronto \textbf{Final} ).
É \textbf{relevante} registrar a \textbf{descrição} da representação social das personagens e o destaque de \textbf{tópicos} relacionados aos objetos de conhecimento, bem como assuntos que causam interesse de pesquisa pelo estudante. - Delimitação e construção de situações-problema ( cenário de aprendizagem ), orientação e acompanhamento de leituras exploratórias, acesso ao material de apoio, produção de fichamentos, resumos. - Roda de conversa para levantamento de temas \textbf{relevantes}, pelos estudantes, referentes aos campos da manipulação e \textbf{edição} genética. - Visitas ao laboratório de Genética ( INBIO/UFMS ) para conhecer os principais projetos e pesquisas desenvolvidas na área de biotecnologia e \textbf{edição} genética. - \textbf{Sistematização} de um projeto de pesquisa ( tema delimitado, objetivos, justificativa, metodologia, cronograma e fontes ). - Fase 1 da execução da pesquisa pelo estudante e acompanhamento processual pelo professor : revisão de literatura ( pesquisa bibliográfica ). - Fase 2 da execução da pesquisa pelo estudante e acompanhamento processual pelo professor : pesquisa de campo e experimental com embasamento estatístico-matemático. - Produção de um ensaio seguindo as normas da ABNT. - Socialização, avaliação compartilhada ( entre os estudantes e o professor ) e revisão do material produzido. - Comunicação dos ensaios da turma na forma de \textbf{seminário}. - Produção de banner para comunicação na feira de \textbf{ciências} ou \textbf{evento} de culminância da escola e, posteriormente, disponibilizar em mídias da escola ( redes sociais, jornais, blogs, dentre outros ), com o fim de produzir esclarecimento e intervenção no contexto sociocultural. \textbf{Problematização} com o uso do Arco de Maguerez, partindo das seguintes etapas : - \textbf{Análise} semiótica de memes, comentários e \textbf{vídeos} \textbf{postados} em redes sociais, que alimentam o preconceito e o racismo genético.
Novamente, é \textbf{relevante} registrar os \textbf{tópicos} e assuntos que causam interesse de pesquisa pelo estudante. - Identificação de situações relacionadas à discriminação étnica ou racial em âmbito local, no cotidiano dos estudantes e/ou levantamento por meio da aplicação de questionário de pesquisa de campo sobre as principais situações discriminatórias e/ou de desigualdade de oportunidades encontradas na comunidade ou até mesmo no ambiente escolar. - Organização e \textbf{análise} dos dados da pesquisa realizada, para delimitação de situação \textbf{problema}. - Execução da pesquisa bibliográfica com acompanhamento processual do professor para enriquecimento e \textbf{contextualização} dos dados da pesquisa ( \textbf{análise} de dados IPEA ). - Realização de rodas de conversa sobre os principais \textbf{problemas} identificados e delimitação de ação de intervenção. - Construção de mapas \textbf{conceituais} para organizar as aprendizagens e conceitos relacionados ao material \textbf{teórico} e desenvolvimento das pesquisas. - Proposição de intervenção relacionada à situação \textbf{problema}. - Produção de material para socialização entre os estudantes e o professor ( avaliação compartilhada e revisão do material produzido ). - Elaboração de material para comunicação em mídias sociais, \textbf{explicitando} os resultados das pesquisas e compartilhamento de experiências ( \textbf{vídeos}, podcasts e/ou canais de youtube ). - Apresentação do material produzido em \textbf{evento} de culminância da escola e/ou em \textbf{eventos} de relevância ligados a ações sobre discriminação étnica ou racial, com o fim de produzir esclarecimento e intervenção no contexto sociocultural.
Avaliação A avaliação é processual e \textbf{supõe} o desenvolvimento de todas as etapas desta atividade de aprendizagem, de modo a contemplar os eixos estruturantes programados ao longo da unidade curricular ; assim devem ser considerados \textbf{tanto} os ensaios e banners utilizados para apresentação em \textbf{eventos} de culminância na escola como as produções multimodais, em especial a produção de gêneros \textbf{digitais}.
Além da avaliação do professor, \textbf{sugere} -se que os colegas também possam avaliar, de forma colaborativa, o material produzido pela sua turma.
Ao avaliar, o professor deve \textbf{verificar} se as produções : a ) atendem ao tema delimitado ; b ) expressam de forma adequada as informações e a \textbf{contextualização} ; c ) apresentam justificativas e \textbf{argumentos} que sustentam a conclusão ; d ) pautam informações \textbf{pertinentes} e diversificadas ; e ) têm caráter autoral, ou seja, que não sejam cópias ( plágios ).
Observações A arquitetura desta unidade curricular envolve desafios pedagógicos específicos, perante os quais estudantes e professor precisam ter em conta, a princípio, os seguintes aspectos : 1.
Esta unidade trata de questões filosófico-científicas \textbf{complexas} sob a perspectiva de uma \textbf{abordagem} \textbf{multidisciplinar} ( \textbf{biologia}, direito, \textbf{filosofia} dentre outras ), na medida em que tematiza \textbf{problemas}, controvérsias e soluções postos nos campos da bioética e biotecnologias.
Em \textbf{tese}, essa \textbf{abordagem} \textbf{induz} uma experiência de aprendizagem multifocal que \textbf{implica} o efetivo comprometimento por parte dos estudantes e professor nas atividades de estudos, pesquisa e produção autoral, sob pena de margear apenas divulgações panfletárias da \textbf{ciência}. 2.
Em \textbf{razão} da complexidade dos temas abordados, recomenda -se que, consensualmente, professor e estudantes façam adequações e/ou delimitação dos objetivos da unidade e objetos de conhecimento que \textbf{julgarem} \textbf{pertinentes}.
Isso pode potencializar, estrategicamente, níveis de \textbf{aprofundamento} da educação científica e o protagonismo dos estudantes no processo de construção do conhecimento. \textbf{Supondo} isso, por exemplo, é possível selecionar e \textbf{sistematizar} diferentes \textbf{focos} de \textbf{abordagem} a serem trabalhados nesta unidade. 3.
As sugestões didáticas apresentadas, estudo de caso e \textbf{problematização} ( Arco de Maguerez ) procuram enlaçar, no conjunto das atividades de construção do conhecimento, os objetos de conhecimento, as técnicas de metodologia científica e as estratégias didáticas.
Nesse sentido, é importante destacar que, embora os eixos estruturantes estejam delimitados em diferentes \textbf{organizadores}, pressupõe -se o desenvolvimento de modo integrado das ações propostas, para dar sentido ao desenvolvimento das habilidades, sem privilegiar um determinado eixo em detrimento de outro. 4.
Os textos e materiais de apoio indicados, bem como as produções, para avaliação ao longo do desenvolvimento da unidade curricular, oferecem diferentes possibilidades ( textos, livros, \textbf{vídeos}, podcasts ), com \textbf{vistas} a proporcionar acessibilidade para atender estudantes com necessidades especiais.
Nesse sentido, optou -se pela exibição de \textbf{vídeos} do youtube que possibilitam acessibilidade como o uso de legendas e a adequação da velocidade da reprodução. 5.
O item Recursos desta unidade menciona, a título de exemplificação, uma hipotética “ parceria ” - no caso, desejável - com o Laboratório de Genética da Universidade Federal do Mato Grosso do Sul.
Esse exemplo hipotético reforça a orientação no sentido de que a escola efetive parcerias institucionais respeitando as orientações do órgão \textbf{central}, bem como vínculos acadêmicos para auxiliá -la em sua \textbf{tarefa} formativa.
Em uma sociedade interconectada, as parcerias da escola com a sociedade civil e a participação em redes de pesquisa e inovação científico-culturais, programadas no contexto da educação, podem contribuir para inserir os estudantes no âmbito da comunidade científica e, a partir disso, subsidiá -los intelectualmente para uma prática de intervenção sociocultural.

% matched lemmas: abordagem, análise, aprofundamento, argumento, biologia, central, ciência, complexo, conceitual, contextualização, crítico, debate, descrição, digital, edição, evento, explicitar, filosofia, final, foco, implicar, imprimir, induzir, julgar, moderno, multidisciplinar, organizador, pertinente, postar, problema, problematização, razão, relevante, seminário, sistematizar, sistematização, sugerir, supor, tanto, tarefa, teoria, tese, teórico, tópico, verificar, vista, vídeo
\end{textsample}
